% Figure: log(N/N0) vs sterilisation time, showing first-order kill
% kinetics and the 6-log and 12-log reduction targets.
\begin{figure}[htbp]
\centering
\begin{tikzpicture}[
  axis/.style={->, thick},
  kill/.style={engblue, very thick},
  target/.style={engred, densely dashed, thick},
  ref/.style={engblack!60, densely dashed},
  lab/.style={font=\small, align=left},
]

\draw[axis] (0, 0) -- (9.5, 0) node[right, lab] {time (multiples of $D$)};
\draw[axis] (0, -6.5) -- (0, 0.5) node[above, lab] {$\log_{10}(N/N_0)$};

% First-order kill: log(N/N0) = -t/D, so slope is -1 per D unit
\draw[kill] (0, 0) -- (9, -9);

% Tick marks on x
\foreach \x in {0,1,...,9}
  {\draw[engblack!60] (\x, 0) -- (\x, 0.1); \node[lab, anchor=south, font=\scriptsize] at (\x, 0.1) {\x};}
% Tick marks on y
\foreach \y in {0,-1,-2,-3,-4,-5,-6}
  {\draw[engblack!60] (0, \y) -- (-0.1, \y); \node[lab, anchor=east, font=\scriptsize] at (-0.1, \y) {\y};}

% 6-log target
\draw[target] (0, -6) -- (6, -6) -- (6, 0);
\node[lab, engred, anchor=south] at (6, 0.0) {$t_{6D}$};
\node[lab, engred, anchor=east] at (0, -6) {6-log};

% 3-log target reference
\draw[ref] (0, -3) -- (3, -3) -- (3, 0);
\node[lab, engblack!70, anchor=south, font=\scriptsize] at (3, 0.0) {$t_{3D}$};

% Annotation box
\node[lab, engblue, anchor=west, draw=engblue!50, fill=engblue!5,
      rounded corners, inner sep=4pt]
  at (5.0, -2.0)
  {$N(t) = N_0 \, 10^{-t/D}$\\
   $t_{nD} = n D$\\
   $D = 2.1\,\text{min}$ (B.\ subtilis spores,\\
   moist heat, $121\,\degC$)};

\end{tikzpicture}
\caption{First-order microbial kill kinetics on log scale, expressed
in multiples of the decimal reduction time $D$. The kill curve is a
straight line with slope $-1$ per $D$ unit, so an $n$-log reduction
requires $t = n D$. A six-log reduction target (red dashed) is the
typical sterile-filtration design goal for sub-micron bacteria; a
twelve-log target is used for terminal heat sterilisation of finished
parenteral products. The $D$ value for \emph{Bacillus subtilis}
spores under saturated steam at $121\,\degC$ is about $2.1$
minutes; the same organism's $D$ rises sharply at lower temperatures
and falls at higher temperatures, following an Arrhenius dependence
captured by the $z$ value.}
\label{fig:vol06ch10:sterile-log-reduction}
\end{figure}
